\section*{F.1 Original Project Proposal Summary}

The dissertation was originally proposed as a systematic audit of AI responses to a standardised battery of mental health queries, varying the implied geographic location of the user across twenty cities spanning different income levels, healthcare infrastructure, and cultural contexts, with responses evaluated across four dimensions: crisis-line accuracy benchmarked against the WHO Mental Health Atlas ground truth, service referral type, cultural assumption coding, and response specificity. A mixed-methods approach combining automatic metrics across all responses with structured human coding on a stratified subset, with inter-rater reliability testing, was proposed throughout.

\section*{F.2 Deviations from the Original Proposal, and Rationale}

Three material deviations from the original proposal are recorded here, consistent with this dissertation's overall commitment to documenting rather than silently absorbing changes to the original design.

\textbf{Model selection.} The original proposal envisaged a different, unspecified set of large language models. The seven models actually used (Section 3.2.3) were selected instead, based on API accessibility and a mix of open-weight and commercial models providing broader coverage of the current LLM landscape than the originally envisaged set.

\textbf{Reduction from four evaluation dimensions to three outcome families.} Following supervisory feedback that the original four-dimension design risked accumulating many weak, loosely connected indicators rather than testing a small number of well-specified hypotheses, the design was reduced to three outcome families -- actionability, localisation, and support orientation -- each explicitly connected to a single pre-registered hypothesis, geographic predictor, and statistical model (Section 3.3--3.5). Crisis-line accuracy, originally proposed as a fourth dimension, is instead addressed as an exploratory finding (the crisis-contact fabrication analysis, Section 4.7) rather than as a fourth confirmatory hypothesis, since it emerged from auditing the data after the confirmatory design was finalised rather than from the pre-registered plan.

\textbf{Crisis-service reference data.} The original proposal specified benchmarking against the WHO Mental Health Atlas as ground truth. In practice, a proxy reference constructed from public crisis-line directories (IASP, Befrienders Worldwide, Find A Helpline) was used, with five of twenty countries directly re-verified against these and, where available, WHO sources during this dissertation (Appendix C.3). Completing this verification against the originally specified WHO Mental Health Atlas for all twenty countries is identified as a priority direction for future work (Section 6.3).

\section*{F.3 Supervisory Feedback Incorporated}

Feedback received during supervision meetings and incorporated into the final design included: aggregating observations at city level before introducing income or regional classifications as variables, to avoid erasing city-level heterogeneity; connecting each hypothesis explicitly to the specific feature used to test it; requiring precision, recall, and F1 to be reported for every automated feature, with weak features revised or discarded rather than retained; retaining response- and city-level data throughout rather than aggregating away city, model, or prompt identifiers; and producing a small number of core figures and tables, each directly answering one hypothesis, rather than a large number of exploratory visualisations without a clear analytical hierarchy.
